We evaluate and compare the Neural SDDDE on several dynamical systems listed below coming from biology and population dynamics. We show that Neural SDDDE outperforms a variety of continuous-depth models and demonstrates its capabilities in simulating state-dependent delayed systems.


\subsection{Dynamical systems}\label{ssec:system_exp}
We pick different dynamical systems that exhibit specific delays since it is our main focus. Some systems are taken from \cite{neural_laplace} for comparison.

\paragraph{One constant time-delay system} We consider the delayed Lotka-Volterra from \cite{lk_dde}, that simulates population dynamics:
%
\begin{equation}
\begin{aligned}
&\dot{x}(t) = \frac{1}{2} x(t) ( 1  - y(t-\tau)) \\
&\dot{y}(t) = -\frac{1}{2} y(t)( 1  - x(t-\tau)). \\
\end{aligned}
\end{equation}
%
The delay is fixed to $\tau=0.2$, we integrate in the time range $[0,30]$ and define the constant history function $\phi(t) = (x_0, y_0)$ $\forall t<0$ with $x_0$ and $y_0$ uniformly sampled from $[0.1,2]$.

\paragraph{Chaotic system with one constant time delay} We choose a chaotic setting of the Mackey-Glass system \citep{mackey_glass} that mimics pathological behavior in certain biological contexts:
%
\begin{equation}
\begin{aligned}
&\dot{x}(t) = \beta \frac{x(t-\tau)}{1 + x(t-\tau)^n} - \gamma x(t),
\end{aligned}
\end{equation}
%
with $\beta=0.25$, $n=10$ and $\gamma=0.1$. The delay is fixed to $\tau= 10$. We integrate in the time range $[0,80]$ and define the constant history function  $\phi(t) = x_0$ where $x_0$ is uniformly sampled from $[0.1,1.0]$.

\paragraph{Time-dependent delay system} Here, the delay is time dependent:
%
\begin{equation}
\begin{aligned}
&\dot{x}(t) = x(t) (1 - x(t-\tau(t))),
\end{aligned}
\end{equation}
%
with $\tau(t) = 2 + \sin(t)$. We integrate in the time range $[0,40]$ and define the constant history function $\phi(t) = x_0$, where $x_0$ is uniformly sampled from $[0.1,2.0]$.

\paragraph{State-dependent time-delay system} In this example, we consider the 1-D state-dependent Mackey Glass system from \citet{state_dependent_mg} with a state-dependent delay:
%
\begin{equation}
\dot{x}(t) = -\alpha(t) x(t) + \beta(t) \frac{ x^2(t - \tau(x))}{1 + x^2(t - \tau(x))} + \gamma(t)
\end{equation}
\begin{equation}
\begin{aligned}
\text{with} \quad \alpha(t) & = 4 + \sin(t) + \sin(\sqrt{2t}) + \frac{1}{1 + t^2} \\
 \beta(t) & = \gamma(t) = \sin(t) + \sin(\sqrt{2t}) + \frac{1}{1 + t^2}
\end{aligned}
\end{equation}
%
and $\tau(x) = \frac{1}{2} \cos{x(t)}$. The model is defined on the time range $[0,10]$ and the constant history function is $\phi(t) = x_0$ with $x_0$ uniformly sampled from $[0.1,1]$.

\paragraph{Delayed Diffusion Equation} Finally, we choose the delayed PDE taken from \citet{arino2009delay}. Such dynamics can for example model single species growth in a food-limited environment.
%
\begin{equation}
\begin{aligned}
&\frac{\partial u}{\partial t}(x,t) = D \frac{\partial^2 u}{\partial x^2}(x,t) + r u(x,t) \left( 1 - u(x,t-\tau) \right),
\end{aligned}
\end{equation}
%
where $D=0.01, r=0.9$ and $\tau=2$. We integrate in the time range $[0,4]$, the spatial domain is $\mathcal{D}_x = [0,1]$ with periodic boundary conditions and define the history function $\phi(x, t) = a \sin{x} e^{-0.01t} $ where $a$ is uniformly sampled from $[0.1,4.0]$. The spatial mesh is created with $dx=0.01$.

\paragraph{Data generation} For all the systems listed in this section, data generation information is gathered in ~\ref{ap:data_generation} along with information of the step history function generation in \ref{ap:step_history_function}.

\begin{table*}[t]
\begin{center}
\begin{small}
\begin{sc}
% \resizebox{14cm}{!}{
\begin{tabular}{l|cccr}
% \toprule
   & NODE &  ANODE & Neural Laplace & Neural SDDDE \\
% \midrule
\hline
Lokta-Volterra & $.00288 \pm .00148$ & $.0132 \pm .0081$ & $.275 \pm .178$ & $ \mathbf{.00265 \pm .00074}$ \\
Mackey-Glass & $.668 \pm 1.22$ &  $ .00355 \pm .00383$& $.0123 \pm .0093$ & $\mathbf{.000159 \pm .000189}$ \\
Time Dependent DDE &  $1.01 \pm .435$ & $.00729 \pm .00235$ & $.0144 \pm .0037$ & $\mathbf{.00148 \pm  .00087}$ \\
State-Dependent DDE & $.0363 \pm .0015$ & $.00107 \pm .00076$ & $.00379 \pm .00110$ & $\mathbf{.000149 \pm .000007}$ \\
Delay Diffusion & $.00117 \pm .00020$ & $.000200 \pm .000052$ & $\mathbf{.0000737 \pm .0000106}$ & $.000300 \pm .000043$ \\
% \bottomrule
\hline
\end{tabular}
% }
\end{sc}
\end{small}
\caption{Test MSE averaged over 5 runs of each experiments with their standard deviation. Best result bolded.}
\label{test_rmse}
\end{center}
\end{table*}


\subsection{Evaluation}
\label{ss:evaluation}
We assess the performance of the models with their ability to predict future states of a given system. The metric used is the mean square error (MSE) in all cases. Neural Laplace predicts only after a burn-in time
since a part of the observed trajectory is used to learn a latent initial condition vector $\mathbf{p}$. Since NODE, ANODE and Neural SDDDE can be seen as initial value problems (IVPs) we produce trajectories from initial conditions and compute the MSE with respect to the whole trajectory.
On each DDE system, to judge the quality of each model, we elaborate two additional experiments alongside with the \textit{test set predictions}. By modifying the history function $\phi(t)$ we change the behavior of the system and hence generate new trajectories. The first experiment puts each model in a \textit{pure extrapolation regime} ; the constant value of the history function $\phi(t)=\mathbf{x}_0$ is sampled outside the range of the training and testing data. This allows to see the models' extrapolation capabilities. The second assessment is more a hybrid approach where the \textit{history function is a step function}:

\begin{equation*}
\begin{aligned}
\phi(t) &= \begin{cases} 
      \mathbf{x}_0 & t\leq t_{\text{jump}} \\
      \mathbf{x}_1 & \text{otherwise}
   \end{cases} \\
t_{\text{jump}} &\sim \mathcal{U}(-\tau_{\text{max}}, 0), \quad \mathbf{x}_0, \mathbf{x}_1 \sim \mathcal{U}(c_0, c_1)
\end{aligned}
\end{equation*}

where $t_{\text{jump}}$ is the largest delay in the system and $c_0, c_1$ are system specific randomly sampled values (see \ref{sss:history_function} for more details). Not only can the nature of the history step function change but can also have its domain function outside of the training and test data (extrapolation regime).

As a reminder, to produce outputs, NODE and ANODE need an initial condition, Neural SDDDE the history function and Neural Laplace a portion of the trajectory. To ensure comparison in our experiments, we opt to give Neural Laplace the same information as Neural SDDDE, specifically the history function. This could be seen as a restriction of Neural Laplace since in its authors give 50\% of the trajectory to produce outputs. Therefore, we propose another small experiment: on one dynamical system, Neural SDDDE is compared with Neural Laplace which is given more than just the history function. 

As stated previously, Neural Laplace does not need to define delays thanks to the Laplace transformation compared to Neural SDDDE. This is why we provide one last experiment: for the Mackey-Glass system, several delays are provided to the vector field of Neural SDDDE. The model will have to learn the system's dynamics with the panel of provided delays.

 \begin{figure*}
    \centering
    \begin{subfigure}[b]{0.475\textwidth}
        \centering
        \includegraphics[width=\textwidth]{final_images/experiments/lk/testset_2_samples.png}
          \caption{Lotka Volterra randomly sampled test trajectory plots}\label{subfig:lk_testset}
    \end{subfigure}
    \hfill
    \begin{subfigure}[b]{0.475\textwidth}  
        \centering 
        \includegraphics[width=\textwidth]{final_images/experiments/mg/testset_2_samples.png}
         \caption{Mackey-Glass randomly sampled test trajectory plots}\label{subfig:mg_testset}
    \end{subfigure}
    \vskip\baselineskip
    \begin{subfigure}[b]{0.475\textwidth}   
        \centering 
        \includegraphics[width=\textwidth]{final_images/experiments/time/trainset_2_samples.png}
         \caption{Time Dependent DDE randomly sampled test trajectory plots }\label{subfig:time_testset}
    \end{subfigure}
    \hfill
    \begin{subfigure}[b]{0.475\textwidth}   
        \centering 
        \includegraphics[width=\textwidth]{final_images/experiments/state/trainset_2_samples.png}
         \caption{State Dependent DDE randomly sampled test trajectory plots }\label{subfig:state_testset}
    \end{subfigure}
    \vskip\baselineskip
    \begin{subfigure}[b]{1.0\textwidth}  
        \centering 
        % \includegraphics[width=\textwidth, trim={0cm 5cm 0 12cm},clip]{final_images/experiments/diffusion/testset.png}
         \includegraphics[width=\textwidth, trim={0cm 2.5cm 0 8cm},clip]{final_images/experiments/diffusion/testset2_other.png}
         \caption{Diffusion Delay PDE  randomly sampled from the testset}\label{subfig:diffusion_testset}
    \end{subfigure}
    \vskip\baselineskip
    \begin{subfigure}[b]{1.0\textwidth} 
        \centering 
        % \includegraphics[width=\textwidth, trim={0cm 5cm 0 12cm},clip]{final_images/experiments/diffusion/testset_diff.png}
        \includegraphics[width=\textwidth, trim={0cm 2.5cm 0 8cm},clip]{final_images/experiments/diffusion/testset2_diff_other.png}
         \caption{Absolute error of Diffusion Delay PDE randomly sampled from the testset}\label{subfig:diffusion_diff_testset}
    \end{subfigure}
    \caption{Randomly test trajectory plots of Lotka-Volterra, Mackey-Glass, Time, State Dependent DDEs and Diffusion Delay PDE
    % \LMcomment{The colormap of the last row (error) is poorly chosen since i) the truth appears like it has the largest error; 2) it is visually unclear which methods are performing better. I would invert the colormap (white is toward zero) and possibly restrict the interval to something like [1e-4, 1e-1] or so. Same comment applies to Fig. 2 below.}
    }
    \label{fig:all_testset}
    \end{figure*}


\subsection{Results}

Test errors for each experiments are reported in Table~\ref{test_rmse}. Complementary information are included in the Appendix \ref{sec:appendix} for what concerns the implementation and the training process; training curves are plotted in ~\ref{ap:training_curves}, model and training hyperparameters are given in \ref{ap:model_hyper}.

\paragraph{Testset prediction} Neural SDDDE almost consistently outperforms each model over all DDE systems mentioned in section \ref{sec:experiments} and plots are provided with Figure \ref{fig:all_testset}. Neural Laplace seems to suffer from the Runge phenomenon where its most evident example is the State-Dependent DDE (Figure \ref{subfig:state_testset}) but also in the Mackey-Glass DDE (Figure \ref{subfig:mg_testset}). This is most likely due to the query algorithm that does not span enough the Laplace domain. Not surprisingly, Neural ODE is the most limited model: all over predicts only the dynamics' mean trajectory of the dynamical systems under consideration except for the Lotka Volterra dataset where it reproduces the correct dynamics (Figure \ref{subfig:lk_testset}). ANODE gives satisfactory results except in the Time Dependent DDE (\ref{subfig:time_testset}). Finally, Neural SDDDE out of all the models provide the best outcome. Regarding the Diffusion Delay PDE, the PDE's evolution is predicted with an absolute error going as high as $10^{-2}$ with all models as seen in Figure \ref{subfig:diffusion_testset}. The absolute error is also displayed in Figure \ref{subfig:diffusion_diff_testset}: it depicts the discrepancy between models. Neural Laplace yields more errors for a given time steps across all the spatial domain compared to IVP models that have errors more localized in specific spatial regions.

\begin{figure*}
    \centering
    \begin{subfigure}[b]{0.475\textwidth}
        \centering
        \includegraphics[width=\textwidth]{final_images/experiments/lk/extrapolate_2_samples.png}
          \caption{Lotka Volterra randomly sampled extrapolated trajectory plots}\label{subfig:lk_extrapolate}
    \end{subfigure}
    \hfill
    \begin{subfigure}[b]{0.475\textwidth}  
        \centering 
        \includegraphics[width=\textwidth]{final_images/experiments/mg/extrapolate_2_samples.png}
         \caption{Mackey-Glass randomly sampled extrapolated trajectory plots}\label{subfig:mg_extrapolate}
    \end{subfigure}
    \vskip\baselineskip
    \begin{subfigure}[b]{0.475\textwidth}   
        \centering 
        \includegraphics[width=\textwidth]{final_images/experiments/time/extrapolate_2_samples.png}
         \caption{Time Dependent DDE randomly sampled extrapolated trajectory plots }\label{subfig:time_extrapolate}
    \end{subfigure}
    \hfill
    \begin{subfigure}[b]{0.475\textwidth}   
        \centering 
        \includegraphics[width=\textwidth]{final_images/experiments/state/extrapolate_2_samples.png}
         \caption{State Dependent DDE randomly sampled extrapolated trajectory plots }\label{subfig:state_extrapolate}
    \end{subfigure}
      \vskip\baselineskip
    \begin{subfigure}[b]{1.0\textwidth}  
        \centering 
        % \includegraphics[width=\textwidth, trim={0cm 5cm 0 12cm},clip]{final_images/experiments/diffusion/extrapolate.png}
        \includegraphics[width=\textwidth, trim={0cm 2.5cm 0 8cm},clip]{final_images/experiments/diffusion/extrapolate_other4.png}
         \caption{Diffusion Delay PDE  randomly sampled from the extrapolated testset}\label{subfig:diffusion_extrapolate}
    \end{subfigure}
    \vskip\baselineskip
    \begin{subfigure}[b]{1.0\textwidth} 
        \centering 
        % \includegraphics[width=\textwidth, trim={0cm 5cm 0 12cm},clip]{final_images/experiments/diffusion/extrapolate_diff.png}
         \includegraphics[width=\textwidth, trim={0cm 2.5cm 0 8cm},clip]{final_images/experiments/diffusion/extrapolate_diff_other4.png}
         \caption{Absolute error of Diffusion Delay PDE randomly sampled from the extrapolated testset}\label{subfig:diffusion_diff_extrapolate}
    \end{subfigure}
    \caption{Randomly extrapolated trajectory plots of Lotka-Volterra, Mackey-Glass, Time, State Dependent DDEs and Diffusion Delay PDE}
\label{fig:all_extrapolate}
    \end{figure*}





\paragraph{Extrapolation regime prediction} This experiment really challenges model generalization capabilities. Overall, on certain datasets, some models can extrapolate with new constant history functions that are not too far out from the function domain of history functions used during training; more in details,  trajectories were generated with $\phi(t) = x_0\in [a,b]$: some models are able to exhibit adequate predictions for history functions that have a value near the bounds of $[a,b]$. For example, for the Mackey-Glass system (Figure \ref{subfig:mg_extrapolate}), Neural SDDDE predictions are good up to values of $x_0 \approx 1.2$. For values higher than  $x_0 \approx 1.4$ (during training $\phi(t) = x_0 \in [0.1,1] $) overall variations are captured but the amplitudes are off. On the same Figure \ref{subfig:mg_extrapolate}, Neural Laplace yields anti-phase results along with what seems to be Runge phenomenon. For the Time Dependent system (Figure \ref{subfig:time_extrapolate}), Neural Laplace yields better results: trajectories possess the overall dynamical trend with some amplitude discrepancy. Regarding performances of NODE and ANODE, results are quite heterogeneous, most of them fail apart for some exceptions discussed below. For the State-Dependent DDE (Figure \ref{subfig:state_extrapolate}), ANODE produces satisfactory trajectories. This can be explained because the system at hand has almost periodic solutions thus easier to learn. For the Lotka Volterra dataset (Figure \ref{subfig:lk_extrapolate}), IVP models (ANODE, NODE, Neural SDDDE) across the board yield realistic results even though the system's dynamics is not captured perfectly. For the Diffusion Delay PDE, overfitting is observed for Neural Laplace. Out of all the IVP models, Neural SDDDE predicts the best possible outcome compared to NODE and ANODE. 

\begin{figure*}[t]
    \centering
    \begin{subfigure}[b]{0.475\textwidth}
        \centering
        \includegraphics[width=\textwidth]{final_images/experiments/lk/other_history_2_samples.png}
          \caption{Lotka Volterra randomly sampled from history step function}\label{subfig:lk_other_history}
    \end{subfigure}
    \hfill
    \begin{subfigure}[b]{0.475\textwidth}  
        \centering 
        \includegraphics[width=\textwidth]{final_images/experiments/mg/other_history_2_samples.png}
         \caption{Mackey-Glass randomly sampled from history step function }\label{subfig:mg_other_history}
    \end{subfigure}
    \vskip\baselineskip
    \begin{subfigure}[b]{0.475\textwidth}   
        \centering 
        \includegraphics[width=\textwidth]{final_images/experiments/time/other_history_2_samples.png}
         \caption{Time Dependent DDE randomly sampled from history step function  }\label{subfig:time_other_history}
    \end{subfigure}
    \hfill
    \begin{subfigure}[b]{0.475\textwidth}   
        \centering 
        \includegraphics[width=\textwidth]{final_images/experiments/state/other_history_2_samples.png}
         \caption{State Dependent DDE randomly sampled from history step function  }\label{subfig:state_other_history}
    \end{subfigure}
    \caption{Randomly sampled trajectory plots from other history step function experiment of Lotka-Volterra, Mackey-Glass, Time, State Dependent DDEs and Diffusion Delay PDE}
  \label{fig:all_other_history}
\end{figure*}

\paragraph{Step history function prediction} This third experiment also demonstrates how the modification of the history function leads to changes in the transient regime and impacts later dynamics. For example, the plots in Figure~\ref{subfig:mg_other_history} exhibits a monotonic change at $t \approx 10$ (it appeared because we added a new discontinuity jump in the history function) that never appeared beforehand and Neural SDDDE captures it to some extent. Neural Laplace fails to generate adequate trajectories apart from the State-Dependent DDE already discussed in the previous paragraph. Regarding the other IVP models exhibit the same performance as in the previous paragraph. By studying the effect of such a new history function on the Diffusion Delay PDE, we saw that the system's dynamics is not changed substantially, therefore, we decided to omit this system's comparison.

\begin{table*}[t]
\begin{center}
\begin{small}
\begin{sc}
\begin{tabular}{l|cccr}
  & NODE &  ANODE & Neural Laplace & Neural SDDDE \\
\hline
 $\alpha=0$ &  $1.01 \pm .435$ & $.00729 \pm .00235$ & $.0144 \pm .00368$ & $\mathbf{.00148 \pm  .000872}$ \\
 $\alpha=0.02$ & $.720 \pm .00254$ & $.0128 \pm .002377$ & $.00881 \pm .00254$ & $\mathbf{.000906 \pm .000441}$\\
 $\alpha=0.05$ &  $4.032 \pm  4.225$ & $.03655 \pm .0349$ & $.00977 \pm .00146$ & $\mathbf{.00250 \pm .000951}$\\
 $\alpha=0.1$ & $1.597 \pm 1.100$ & $.0223 \pm .00634$ & $.0154 \pm .00501$ & $\mathbf{.0121 \pm .00534}$\\
 $\alpha=0.2$ & $1.02 \pm .282$ & $.0321 \pm .00319$ & $.0273 \pm .00704$ & $\mathbf{.0186 \pm .00524}$\\
\hline
\end{tabular}
\end{sc}
\end{small}
\caption{Test MSE with the noiseless data averaged over 5 runs of each Time Dependent DDE noise experiments with their standard deviation. Best result bolded.}
\label{noiseless_mse}
\end{center}
\end{table*}

\paragraph{Increasing trajectory fed for Neural Laplace} In this trial, instead of only giving the history function to Neural Laplace we provide 50\% of the trajectory to construct its latent initial condition representation vector $\mathbf{p}$ during training as in its original paper. Thus, the first half of the trajectory is injected in Neural Laplace to predict the second half. We choose to train the model on the Time Dependent DDE along with the same training procedure (see  \ref{ap:model_hyper}). We provide the training loss for this new setup in Figure \ref{fig:50_neural_laplace_training_curve} where Neural Laplace achieves the lowest training data. Provided the test MSE of Table \ref{test_rmse}, we choose to only compare the test MSE of Neural SDDDE and Neural Laplace in Table \ref{50_laplace_mse}. By comparing Figure \ref{fig:50_neural_laplace_comparison} and \ref{subfig:time_testset} one can clearly note that the Runge phenomenon are almost absent and predictions are almost as good as Neural SDDDE. This confirms that, in general, Neural Laplace needs more than the history function in order to correctly simulate DDEs. 

\begin{figure}[H]
\centering
\includegraphics[width=0.5\textwidth,height=0.2\textheight, keepaspectratio]{final_images/experiments/time/trainset_50percent_laplace.png}\caption{Time-dependent DDE randomly sampled testset trajectories where 50\% of data is fed to Neural Laplace}
\label{fig:50_neural_laplace_comparison}
\end{figure}

\begin{table}[H]
\begin{center}
\begin{small}
\begin{sc}
\begin{tabular}{l|cr}
& Neural Laplace & Neural SDDDE \\
\hline
Test MSE & $\mathbf{.000543 \pm .000235}$ & $.00148 \pm  .00087$ \\
\hline
\end{tabular}
\end{sc}
\end{small}
\caption{Time Dependent test MSE averaged over 5 runs of each experiments with their standard deviation. Best result bolded.}
\label{50_laplace_mse}
\end{center}
\end{table}


\paragraph{Noise analysis} We conduct a noise study on one of the datasets, the Time Dependent DDE system. Each data point is added Gaussian noise that is scaled with a certain factor $\alpha$ of the trajectory's variance. The model is then trained with this noisy data and evaluated on the noiseless testset. In our experiment, we selected 4 scaling factors $\alpha$: $0.02, 0.05, 0.1$ and $0.2$. Results in Table \ref{noiseless_mse} show that our model is robust to noisy data and almost consistently outperforms other models. Additionally, results from Table \ref{noiseless_mse} show that adding a small amount of noise (here $\alpha = 0.02$) makes the learning process more robust, a common result in Machine Learning \citep{noise_robust1,noise_robust2}. 

\paragraph{Giving several delays to choose from} We now consider the case where Neural SDDDE is instantiated with a certain number of delays but only a subset of them correspond to the delayed terms of the system considered in hand. Figure \ref{fig:mg_several_delays} demonstrates that Neural SDDDE is able to provide accurate dynamics for the Mackey-Glass DDE where the delays provided were  $\left\{5.0,10.0,15.0\right\}$ (The training procedure is the same as depicted in  \ref{ap:model_hyper}).

\begin{figure}[H]
\centering
\includegraphics[width=0.5\textwidth,height=0.25\textheight, keepaspectratio]{final_images/mg_several_delays.png}
\caption{Randomly sampled test set of Mackey-Glass system where several delays were provided}
\label{fig:mg_several_delays}
\end{figure}